% !TeX root = ../../report.tex
\subsection{Model Sharpness and Generalisation}\label{sec:model-sharpness-generalisation}

A well-documented limitation of deep learning models applied to high-dimensional biomedical data is their tendency to converge to sharp minima in the loss landscape, which may be associated with poor generalisation and unstable learned representations. Li et al.\ \citeyearpar{li2018} demonstrate that neural networks can converge to either sharp or flat minima within the loss surface \citep{li2018}. Sharp minima are characterised by high sensitivity to parameter perturbations, meaning that small changes in weights can lead to large increases in loss. In contrast, flat minima exhibit smoother regions where the loss remains relatively stable under perturbation. Flat minima are generally associated with improved generalisation, while sharp minima may indicate a greater risk of overfitting and reduced performance on unseen data. This phenomenon is also reflected in the generalisation gap, where a large discrepancy between training and validation or test performance may indicate that the model is learning patterns that do not generalise well \citep{keskar2017}.

The relationship between loss-landscape geometry and explanation quality is less established. Dombrowski et al.\ \citeyearpar{dombrowski2019} show that explanation methods can be sensitive to the geometry of a model's decision boundary, with smoother boundaries potentially producing more stable attributions \citep{dombrowski2019}. Springer et al.\ \citeyearpar{springer2024} further investigate how optimisers targeting flat minima can improve learned feature representations \citep{springer2024}. However, whether these improvements translate into measurable gains in the quality of symbolic rules extracted from a model, particularly in the pharmacogenomic domain, remains insufficiently explored. We therefore hypothesise that improving generalisation through landscape-aware regularisation may improve the fidelity, comprehensibility, and stability of extracted rules. This relationship forms a central hypothesis of this project, introduced in Section~\ref{sec:rqs} and addressed as a research gap in Section 3.1.